\chapter{PHÂN TÍCH YÊU CẦU}

\section{Yêu cầu chức năng}

\subsection{Quản lý tài khoản và phân quyền}
\begin{itemize}
    \item Hệ thống cho phép người dùng đăng ký, đăng nhập, cập nhật thông tin cá nhân và quản lý địa chỉ liên hệ.
    \item Hệ thống quản lý vai trò và quyền truy cập của người dùng; một người dùng có thể đảm nhận nhiều vai trò.
    \item Quản trị viên có thể khóa, mở khóa tài khoản và quản lý quyền sử dụng các chức năng của hệ thống.
\end{itemize}

\subsection{Quản lý Film Lab và dịch vụ}
\begin{itemize}
    \item Film Lab có thể quản lý hồ sơ, địa chỉ, thời gian làm việc và trạng thái hoạt động.
    \item Film Lab có thể tạo và cập nhật dịch vụ tráng phim, scan, in ảnh, gói dịch vụ, giá và thời gian xử lý dự kiến.
    \item Người dùng có thể tìm kiếm và so sánh Film Lab theo vị trí, loại dịch vụ, giá, thời gian xử lý và đánh giá.
    \item Quản trị viên có thể xét duyệt Film Lab và quản lý danh mục dịch vụ dùng chung.
\end{itemize}

\subsection{Quản lý đơn hàng và xử lý phim}
\begin{itemize}
    \item Khách hàng có thể tạo đơn hàng, lựa chọn dịch vụ hoặc gói dịch vụ, khai báo thông tin cuộn phim và yêu cầu bổ sung.
    \item Hệ thống ghi nhận giá, số lượng, phí, giảm giá và tổng tiền tại thời điểm xác nhận đơn hàng.
    \item Film Lab có thể tiếp nhận đơn, quản lý từng cuộn phim và cập nhật các công đoạn như tiếp nhận, tráng, rửa, sấy, scan và kiểm tra chất lượng.
    \item Khách hàng có thể theo dõi trạng thái đơn hàng và tiến độ xử lý phim.
    \item Hệ thống lưu lịch sử thay đổi trạng thái, thời gian thực hiện và người phụ trách để phục vụ tra cứu.
\end{itemize}

\subsection{Giao nhận và thanh toán}
\begin{itemize}
    \item Khách hàng có thể tự giao phim hoặc yêu cầu đối tác đến nhận và giao trả phim.
    \item Hệ thống ghi nhận đối tác giao nhận, địa chỉ, thời gian và trạng thái của từng yêu cầu giao nhận.
    \item Hệ thống hỗ trợ ghi nhận giao dịch thanh toán, phương thức thanh toán, số tiền và trạng thái giao dịch.
    \item Người dùng được thông báo khi trạng thái đơn hàng, giao nhận hoặc thanh toán thay đổi.
\end{itemize}

\subsection{Quản lý kho ảnh phim số}
\begin{itemize}
    \item Film Lab có thể tải lên và bàn giao file scan sau khi hoàn tất xử lý.
    \item Khách hàng có thể xem, tải xuống và tổ chức ảnh theo album hoặc thẻ.
    \item Hệ thống lưu địa chỉ file và metadata như loại phim, máy ảnh, ống kính, ngày chụp, định dạng và thông số scan.
    \item Hệ thống hỗ trợ quản lý các phiên bản file khi ảnh được scan lại.
\end{itemize}

\subsection{Quản lý chợ thiết bị}
\begin{itemize}
    \item Người dùng có thể đăng, cập nhật, tìm kiếm và ngừng hiển thị tin mua bán hoặc trao đổi máy ảnh, ống kính, phim và phụ kiện.
    \item Hệ thống quản lý thông tin sản phẩm, hình ảnh, giá, tình trạng và trạng thái tin đăng.
    \item Hệ thống ghi nhận giao dịch giữa người mua và người bán, đồng thời hỗ trợ đánh giá sau khi giao dịch hoàn tất.
\end{itemize}

\subsection{Tri thức và hoạt động cộng đồng}
\begin{itemize}
    \item Chuyên gia và thành viên được cấp quyền có thể đăng bài hướng dẫn, bài đánh giá và nội dung chia sẻ kinh nghiệm.
    \item Người dùng có thể tìm kiếm bài viết, bình luận và tương tác với nội dung cộng đồng.
    \item Hệ thống hỗ trợ tạo sự kiện, workshop hoặc photowalk và quản lý danh sách đăng ký tham gia.
    \item Quản trị viên có thể kiểm duyệt bài viết, bình luận, đánh giá và xử lý nội dung vi phạm.
\end{itemize}

\subsection{Chức năng hỗ trợ bởi AI}
\begin{itemize}
    \item Hệ thống gợi ý Film Lab và loại phim phù hợp dựa trên nhu cầu, vị trí, sở thích và lịch sử sử dụng của người dùng.
    \item Hệ thống hỗ trợ tìm kiếm ngữ nghĩa trong kho kiến thức nhiếp ảnh.
    \item Trợ lý AI có thể giải đáp câu hỏi về nhiếp ảnh phim và đề xuất nội dung liên quan.
    \item Hệ thống có thể ghi nhận kết quả phân tích chất lượng ảnh scan để hỗ trợ Film Lab và khách hàng.
\end{itemize}
